
\documentclass{article}
\usepackage{colortbl}
\usepackage{makecell}
\usepackage{multirow}
\usepackage{supertabular}

\begin{document}

\newcounter{utterance}

\centering \large Interaction Transcript for game `cladder', experiment `full\_v1.5\_default', episode 4010 with Qwen3.5{-}4B.
\vspace{24pt}

{ \footnotesize  \setcounter{utterance}{1}
\setlength{\tabcolsep}{0pt}
\begin{supertabular}{c@{$\;$}|p{.15\linewidth}@{}p{.15\linewidth}p{.15\linewidth}p{.15\linewidth}p{.15\linewidth}p{.15\linewidth}}
    \# & \multicolumn{2}{c}{Player} && \multicolumn{2}{c}{Game Master} \\
    \hline

    \theutterance \stepcounter{utterance}  
    & & & \multicolumn{4}{p{0.6\linewidth}}{
        \cellcolor[rgb]{0.9,0.9,0.9}{
            \makecell[{{p{\linewidth}}}]{
                \texttt{\tiny{[P1$\langle$GM]}}
                \texttt{You are an expert in causal inference. The following question is not a typical commonsense query, but rather a meticulously designed question created by a professor specializing in causal inference, intended to assess the students' mastery of the course content.} \\
\texttt{Imagine a self{-}contained, hypothetical world with only the following conditions, and without any unmentioned factors or causal relationships: Liking spicy food has a direct effect on parents' social status and child's intelligence. Other unobserved factors has a direct effect on parents' social status and child's intelligence. Parents' social status has a direct effect on child's intelligence. Other unobserved factors is unobserved. The overall probability of liking spicy food is 51\%. The probability of not liking spicy food and intelligent child is 18\%. The probability of liking spicy food and intelligent child is 34\%. Is the chance of intelligent child larger when observing liking spicy food?} \\
\texttt{Start your answer with 'Yes' or 'No', followed by additional reasoning or evidence to support your explanation.} \\
            }
        }
    }
     \\ \\

    \theutterance \stepcounter{utterance}  
    & \multicolumn{4}{p{0.6\linewidth}}{
        \cellcolor[rgb]{0.561,0.737,0.561}{
            \makecell[{{p{\linewidth}}}]{
                \texttt{\tiny{[P1$\rangle$GM]}}
                \texttt{Yes} \\
\\ 
\texttt{**Reasoning:**} \\
\\ 
\texttt{To determine whether the chance of an intelligent child is larger when observing that the child likes spicy food, we must compare the conditional probability \$P(\textbackslash{}text\{Intelligent\} \textbackslash{}mid \textbackslash{}text\{Likes Spicy\})\$ with the marginal probability \$P(\textbackslash{}text\{Intelligent\})\$.} \\
\\ 
\texttt{**Step 1: Calculate the marginal probability of an intelligent child, \$P(\textbackslash{}text\{Intelligent\})\$.**} \\
\texttt{We are given the joint probabilities for the two mutually exclusive and exhaustive states regarding spicy food preference:} \\
\texttt{1.  \$P(\textbackslash{}text\{Not Likes Spicy\} \textbackslash{}cap \textbackslash{}text\{Intelligent\}) = 18\textbackslash{}\%\$} \\
\texttt{2.  \$P(\textbackslash{}text\{Likes Spicy\} \textbackslash{}cap \textbackslash{}text\{Intelligent\}) = 34\textbackslash{}\%\$} \\
\\ 
\texttt{Summing these gives the total probability of an intelligent child:} \\
\texttt{\$\$P(\textbackslash{}text\{Intelligent\}) = 18\textbackslash{}\% + 34\textbackslash{}\% = 52\textbackslash{}\%\$\$} \\
\\ 
\texttt{**Step 2: Calculate the conditional probability of an intelligent child given the child likes spicy food, \$P(\textbackslash{}text\{Intelligent\} \textbackslash{}mid \textbackslash{}text\{Likes Spicy\})\$.**} \\
\texttt{Using the definition of conditional probability:} \\
\texttt{\$\$P(\textbackslash{}text\{Intelligent\} \textbackslash{}mid \textbackslash{}text\{Likes Spicy\}) = \textbackslash{}frac\{P(\textbackslash{}text\{Intelligent\} \textbackslash{}cap \textbackslash{}text\{Likes Spicy\})\}\{P(\textbackslash{}text\{Likes Spicy\})\}\$\$} \\
\\ 
\texttt{First, we determine \$P(\textbackslash{}text\{Likes Spicy\})\$. We know the overall probability of liking spicy food is \$51\textbackslash{}\%\$.} \\
\texttt{\$\$P(\textbackslash{}text\{Likes Spicy\}) = 51\textbackslash{}\%\$\$} \\
\\ 
\texttt{Now, substitute the known values:} \\
\texttt{\$\$P(\textbackslash{}text\{Intelligent\} \textbackslash{}mid \textbackslash{}text\{Likes Spicy\}) = \textbackslash{}frac\{34\textbackslash{}\%\}\{51\textbackslash{}\%\} \textbackslash{}approx 66.67\textbackslash{}\%\$\$} \\
\\ 
\texttt{**Step 3: Compare the probabilities.**} \\
\texttt{*   \$P(\textbackslash{}text\{Intelligent\}) = 52\textbackslash{}\%\$} \\
\texttt{*   \$P(\textbackslash{}text\{Intelligent\} \textbackslash{}mid \textbackslash{}text\{Likes Spicy\}) \textbackslash{}approx 66.67\textbackslash{}\%\$} \\
\\ 
\texttt{Since \$66.67\textbackslash{}\% > 52\textbackslash{}\%\$, the probability of the child being intelligent is indeed higher when we observe that the child likes spicy food.} \\
\\ 
\texttt{**Causal Context Note:**} \\
\texttt{While the problem statement describes a complex causal structure involving confounding variables (unobserved factors affecting both parents' status and child's intelligence, and liking spicy food affecting both), the question strictly asks about the statistical association based on the provided probabilities. In this specific hypothetical world, the observed data (\$34\textbackslash{}\%\$ joint vs \$52\textbackslash{}\%\$ marginal) confirms a positive association. The fact that "liking spicy food" is a cause of "} \\
            }
        }
    }
    & & \\ \\

    \theutterance \stepcounter{utterance}  
    & & & \multicolumn{2}{p{0.3\linewidth}}{
        \cellcolor[rgb]{0.9,0.9,0.9}{
            \makecell[{{p{\linewidth}}}]{
                \texttt{\tiny{[GM$|$GM]}}
                \texttt{Yes} \\
            }
        }
    }
    & & \\ \\

    \theutterance \stepcounter{utterance}  
    & & & \multicolumn{2}{p{0.3\linewidth}}{
        \cellcolor[rgb]{0.9,0.9,0.9}{
            \makecell[{{p{\linewidth}}}]{
                \texttt{\tiny{[GM$|$GM]}}
                \texttt{yes} \\
            }
        }
    }
    & & \\ \\

    \theutterance \stepcounter{utterance}  
    & & & \multicolumn{2}{p{0.3\linewidth}}{
        \cellcolor[rgb]{0.9,0.9,0.9}{
            \makecell[{{p{\linewidth}}}]{
                \texttt{\tiny{[GM$|$GM]}}
                \texttt{game\_result = WIN} \\
            }
        }
    }
    & & \\ \\

\end{supertabular}
}

\end{document}
